% Generated draft: requires mathematical and source review.
\documentclass[11pt]{article}
\usepackage[margin=25mm]{geometry}
\usepackage{fontspec}
\setmainfont{Latin Modern Roman}
\usepackage{amsmath,amssymb,mathtools,unicode-math}
\setmathfont{Latin Modern Math}
\usepackage{enumitem,needspace,xurl,hyperref}
\setlength{\parindent}{0pt}
\setlength{\parskip}{5pt}
\setlength{\emergencystretch}{3em}
\newcommand{\N}{\mathbb{N}}
\newcommand{\Z}{\mathbb{Z}}
\newcommand{\Q}{\mathbb{Q}}
\newcommand{\R}{\mathbb{R}}
\newcommand{\C}{\mathbb{C}}
\newcommand{\F}{\mathbb{F}}
\newcommand{\A}{\mathbb{A}}
\newcommand{\PP}{\mathbb{P}}
\newcommand{\E}{\mathbb{E}}
\newcommand{\eps}{\varepsilon}
\newcommand{\dd}{\,\mathrm{d}}
\begin{document}
\section*{UnsolvedMath ID: 20000374}
\textbf{Source problem code:} catalogue entry
\subsection{Definitions and mathematical statement}
Let \(X\) be a smooth geometrically integral variety over a number field \(K\), let \(S\) be a finite set of places, and let \(Z\subset X\) be closed with codimension at least two. Strong approximation outside \(S\) means density of \(X(K)\) in the projection of its adelic space to the restricted product away from \(S\), with the intended nonempty local-solubility convention. If \(X\) satisfies it, must \(U=X\setminus Z\) also satisfy it? The source first says codimension two and later uses at least two. It also asks for Zariski density of \(S\)-integral points, which requires a chosen integral model and nonempty admissible integral adelic conditions.

The source reports the affine-space case known and highlights smooth affine quadrics \(q(x_1,\ldots,x_4)=c\), simply connected linear algebraic groups, high-dimensional fixed-degree hypersurfaces via the circle method, and quadrics in five variables. It asks whether codimension one could be allowed. Harari's comments invoke algebraic simple connectivity, its persistence under codimension-two removal, and the expectation that Brauer groups will not supply the missing obstruction. These comments are source guidance with implicit hypotheses, not a replacement proof of arithmetic purity.

A related question, attributed to Colliot-Th\'el\`ene, asks for simple connectivity independently of the circle method for diagonal hypersurfaces
\[
\sum_{i=1}^n\epsilon_i x_i^{r_i}=c,\qquad
\epsilon_i\in\{1,-1\},\quad r_i\geq1,\quad c\in\Z,
\]
including the alternating-sign examples in sufficiently many variables handled by that method.

The attached aid reports the general arithmetic-purity question still open and the four-/five-variable quadric cases now covered under the relevant strong-approximation noncompactness hypothesis. It supplies no precise theorem references or full field/local hypotheses for those cases.

Its finite-/tail-sieve reduction separates finitely many local avoidance requirements from a large-prime tail. For candidate point sets \(A(B)\), that tail has the form
\[
\lim_{z\to\infty}\limsup_{B\to\infty}
\frac{\#\{x\in A(B):x\bmod\mathfrak p\in\mathcal Z(\kappa(\mathfrak p))
\text{ for some }N\mathfrak p>z\}}{\#A(B)}=0,
\]
with the integral model and counting setup fixed. The aid reports a summable good-prime local-risk estimate of codimension-two order \(O((N\mathfrak p)^{-2})\), but not that this alone proves the needed uniform tail. It also reports the integral Zariski-density consequence of strong approximation under suitable nonempty integral local conditions, and an explicit codimension-one counterexample, whose data are absent here.

For topology its stated result is precise: over \(\C\), take \(a\ne0\), \(n\geq3\), and
\[
X_a=\{\sum_i\epsilon_i x_i^{r_i}=a\}\subset\A^n_{\C}.
\]
If one exponent is one, it reports \(X_a\simeq\A^{n-1}\). Otherwise it reports homotopy type
\[
X_a\simeq\bigvee^{\prod_i(r_i-1)}S^{n-1}.
\]
It further reports that deleting any closed complex subvariety of codimension at least two leaves \(X_a\) simply connected, hence also with trivial geometric \'etale fundamental group. The nonzero-fiber restriction is material: the source's unqualified \(c\in\Z\) also permits a singular zero fiber, outside this theorem. The topology result does not itself prove strong approximation.
\subsection{Short English statement}
Does deleting a codimension-two subset preserve strong approximation, and what follows for integral points? The attached aid reports qualified quadric cases, isolates a large-prime sieve requirement, and answers a related nonzero diagonal-fiber topology question, while retaining the general arithmetic-purity problem as open.
\subsection{Sources}
\begin{itemize}
\item Local UnsolvedMath registry, permanent ID \(20000374\), exact public code \texttt{catalogue entry}; canonical location \texttt{aim-algebraic-number-theory-notes.json}, notes[25].
\item AIM workshop ``Rational and integral points on higher-dimensional varieties,'' Problem 6; supplied URL \url{https://aimath.org/pastworkshops/ratlhigherdimvarproblems.pdf}. Dialogue is attributed to Wittenberg, Tschinkel, Harari, Wooley, Heath--Brown and Colliot-Th\'el\`ene. The supplied verification reports checking the PDF to repair exponents; this drafting step did not consult it. The literature field is empty.
\item Attached research attempt 1, ``Arithmetic purity: solved quadric cases, a sieve criterion, and diagonal topology,'' referring to cited quadric literature without giving bibliographic details in this record.
\end{itemize}
\end{document}
